% Options for packages loaded elsewhere
\PassOptionsToPackage{unicode}{hyperref}
\PassOptionsToPackage{hyphens}{url}
\documentclass[
]{article}
\usepackage{xcolor}
\usepackage[margin=1in]{geometry}
\usepackage{amsmath,amssymb}
\setcounter{secnumdepth}{-\maxdimen} % remove section numbering
\usepackage{iftex}
\ifPDFTeX
  \usepackage[T1]{fontenc}
  \usepackage[utf8]{inputenc}
  \usepackage{textcomp} % provide euro and other symbols
\else % if luatex or xetex
  \usepackage{unicode-math} % this also loads fontspec
  \defaultfontfeatures{Scale=MatchLowercase}
  \defaultfontfeatures[\rmfamily]{Ligatures=TeX,Scale=1}
\fi
\usepackage{lmodern}
\ifPDFTeX\else
  % xetex/luatex font selection
\fi
% Use upquote if available, for straight quotes in verbatim environments
\IfFileExists{upquote.sty}{\usepackage{upquote}}{}
\IfFileExists{microtype.sty}{% use microtype if available
  \usepackage[]{microtype}
  \UseMicrotypeSet[protrusion]{basicmath} % disable protrusion for tt fonts
}{}
\makeatletter
\@ifundefined{KOMAClassName}{% if non-KOMA class
  \IfFileExists{parskip.sty}{%
    \usepackage{parskip}
  }{% else
    \setlength{\parindent}{0pt}
    \setlength{\parskip}{6pt plus 2pt minus 1pt}}
}{% if KOMA class
  \KOMAoptions{parskip=half}}
\makeatother
\usepackage{color}
\usepackage{fancyvrb}
\newcommand{\VerbBar}{|}
\newcommand{\VERB}{\Verb[commandchars=\\\{\}]}
\DefineVerbatimEnvironment{Highlighting}{Verbatim}{commandchars=\\\{\}}
% Add ',fontsize=\small' for more characters per line
\usepackage{framed}
\definecolor{shadecolor}{RGB}{248,248,248}
\newenvironment{Shaded}{\begin{snugshade}}{\end{snugshade}}
\newcommand{\AlertTok}[1]{\textcolor[rgb]{0.94,0.16,0.16}{#1}}
\newcommand{\AnnotationTok}[1]{\textcolor[rgb]{0.56,0.35,0.01}{\textbf{\textit{#1}}}}
\newcommand{\AttributeTok}[1]{\textcolor[rgb]{0.13,0.29,0.53}{#1}}
\newcommand{\BaseNTok}[1]{\textcolor[rgb]{0.00,0.00,0.81}{#1}}
\newcommand{\BuiltInTok}[1]{#1}
\newcommand{\CharTok}[1]{\textcolor[rgb]{0.31,0.60,0.02}{#1}}
\newcommand{\CommentTok}[1]{\textcolor[rgb]{0.56,0.35,0.01}{\textit{#1}}}
\newcommand{\CommentVarTok}[1]{\textcolor[rgb]{0.56,0.35,0.01}{\textbf{\textit{#1}}}}
\newcommand{\ConstantTok}[1]{\textcolor[rgb]{0.56,0.35,0.01}{#1}}
\newcommand{\ControlFlowTok}[1]{\textcolor[rgb]{0.13,0.29,0.53}{\textbf{#1}}}
\newcommand{\DataTypeTok}[1]{\textcolor[rgb]{0.13,0.29,0.53}{#1}}
\newcommand{\DecValTok}[1]{\textcolor[rgb]{0.00,0.00,0.81}{#1}}
\newcommand{\DocumentationTok}[1]{\textcolor[rgb]{0.56,0.35,0.01}{\textbf{\textit{#1}}}}
\newcommand{\ErrorTok}[1]{\textcolor[rgb]{0.64,0.00,0.00}{\textbf{#1}}}
\newcommand{\ExtensionTok}[1]{#1}
\newcommand{\FloatTok}[1]{\textcolor[rgb]{0.00,0.00,0.81}{#1}}
\newcommand{\FunctionTok}[1]{\textcolor[rgb]{0.13,0.29,0.53}{\textbf{#1}}}
\newcommand{\ImportTok}[1]{#1}
\newcommand{\InformationTok}[1]{\textcolor[rgb]{0.56,0.35,0.01}{\textbf{\textit{#1}}}}
\newcommand{\KeywordTok}[1]{\textcolor[rgb]{0.13,0.29,0.53}{\textbf{#1}}}
\newcommand{\NormalTok}[1]{#1}
\newcommand{\OperatorTok}[1]{\textcolor[rgb]{0.81,0.36,0.00}{\textbf{#1}}}
\newcommand{\OtherTok}[1]{\textcolor[rgb]{0.56,0.35,0.01}{#1}}
\newcommand{\PreprocessorTok}[1]{\textcolor[rgb]{0.56,0.35,0.01}{\textit{#1}}}
\newcommand{\RegionMarkerTok}[1]{#1}
\newcommand{\SpecialCharTok}[1]{\textcolor[rgb]{0.81,0.36,0.00}{\textbf{#1}}}
\newcommand{\SpecialStringTok}[1]{\textcolor[rgb]{0.31,0.60,0.02}{#1}}
\newcommand{\StringTok}[1]{\textcolor[rgb]{0.31,0.60,0.02}{#1}}
\newcommand{\VariableTok}[1]{\textcolor[rgb]{0.00,0.00,0.00}{#1}}
\newcommand{\VerbatimStringTok}[1]{\textcolor[rgb]{0.31,0.60,0.02}{#1}}
\newcommand{\WarningTok}[1]{\textcolor[rgb]{0.56,0.35,0.01}{\textbf{\textit{#1}}}}
\usepackage{graphicx}
\makeatletter
\newsavebox\pandoc@box
\newcommand*\pandocbounded[1]{% scales image to fit in text height/width
  \sbox\pandoc@box{#1}%
  \Gscale@div\@tempa{\textheight}{\dimexpr\ht\pandoc@box+\dp\pandoc@box\relax}%
  \Gscale@div\@tempb{\linewidth}{\wd\pandoc@box}%
  \ifdim\@tempb\p@<\@tempa\p@\let\@tempa\@tempb\fi% select the smaller of both
  \ifdim\@tempa\p@<\p@\scalebox{\@tempa}{\usebox\pandoc@box}%
  \else\usebox{\pandoc@box}%
  \fi%
}
% Set default figure placement to htbp
\def\fps@figure{htbp}
\makeatother
\setlength{\emergencystretch}{3em} % prevent overfull lines
\providecommand{\tightlist}{%
  \setlength{\itemsep}{0pt}\setlength{\parskip}{0pt}}
\usepackage{bookmark}
\IfFileExists{xurl.sty}{\usepackage{xurl}}{} % add URL line breaks if available
\urlstyle{same}
\hypersetup{
  pdftitle={simu},
  hidelinks,
  pdfcreator={LaTeX via pandoc}}

\title{simu}
\author{}
\date{\vspace{-2.5em}2026-05-26}

\begin{document}
\maketitle

Run simulation 2 for Toscano \& McMurray (2010).

Let \(x\) be the observed cue value for a given cue dimension
(e.g.~VOT);

Let \(i \in \{1, \dots, K\}\) index the K Gaussian components that
together form the learner's internal category representation of a given
cue dimension;

Let \(\phi_i\) be the mixing weight, frequency of occurrence) of the
i-th Gaussian;

Let \(\mu_i\) be its mean, representing the central or prototypical cue
value of the component;

Let \(\sigma_i\) be its standard deviation, representing the component's
spread or uncertainty along the cue dimension;

Let \(G_i(x)\) be the likelihood that the i-th Gaussian generated the
observation \(x\) given that cue dimension,

Let \(\eta_\phi\), \(\eta_\mu\), and \(\eta_\sigma\) be the learning
rates for the trial-by-trial updates of \(\Delta\phi_i\),
\(\Delta\mu_i\), and \(\Delta\sigma_i\) to the frequency, mean, and
standard deviation of each Gaussian, respectively.

\[G_i(x) = \phi_i \frac{1}{\sqrt{2\pi\sigma_i^2}} \exp\left(-\frac{(x-\mu_i)^2}{2\sigma_i^2}\right) \tag{4} \]
\[ M(x) = \sum_{i}^{K} G_i(x)  \tag{5} \]

Learning rule:

\begin{itemize}
\tightlist
\item
  A crucial innovation from this model is the use of a winner-take-all
  update rule such that only one Gaussian updates its value on any given
  trial. This allows the model to suppress unneeded categories and
  arrive at the correct solution. Without it, the model does not
  determine the correct number of categories (see McMurray et al.,
  2009a).
\end{itemize}

\[\Delta\phi_i = \eta_\phi \frac{G_i(x)}{M(x)} \tag{A2} \]

\[\Delta\mu_i = \eta_\mu \left(\frac{G_i(x)}{M(x)}\right) \frac{(x-\mu_i)}{\sigma_i^2} \tag{A3} \]

\[\Delta\sigma_i = \eta_\sigma \left(\frac{G_i(x)}{M(x)}\right) \left(\sigma_i^{-3}(x-\mu_i)^2 - \sigma_i^{-1}\right) \tag{A4} \]

\section{1 generate data according to table
1}\label{generate-data-according-to-table-1}

\begin{Shaded}
\begin{Highlighting}[]
\FunctionTok{library}\NormalTok{(ggplot2)}

\CommentTok{\#synthesize data according to table 1}
\NormalTok{d.vot.voiced }\OtherTok{\textless{}{-}} \FunctionTok{rnorm}\NormalTok{(}\AttributeTok{n =} \DecValTok{20000}\NormalTok{, }\AttributeTok{mean =} \DecValTok{0}\NormalTok{, }\AttributeTok{sd =} \DecValTok{5}\NormalTok{)}
\NormalTok{d.vl.voiced }\OtherTok{\textless{}{-}} \FunctionTok{rnorm}\NormalTok{(}\AttributeTok{n =} \DecValTok{20000}\NormalTok{, }\AttributeTok{mean =} \DecValTok{188}\NormalTok{, }\AttributeTok{sd =} \DecValTok{45}\NormalTok{)}
\NormalTok{d.third.voiced }\OtherTok{\textless{}{-}} \FunctionTok{rnorm}\NormalTok{(}\AttributeTok{n =} \DecValTok{20000}\NormalTok{, }\AttributeTok{mean =} \DecValTok{260}\NormalTok{, }\AttributeTok{sd =} \DecValTok{10}\NormalTok{)}

\NormalTok{d.vot.voiceless }\OtherTok{\textless{}{-}} \FunctionTok{rnorm}\NormalTok{(}\AttributeTok{n =} \DecValTok{20000}\NormalTok{, }\AttributeTok{mean =} \DecValTok{50}\NormalTok{, }\AttributeTok{sd =} \DecValTok{10}\NormalTok{)}
\NormalTok{d.vl.voiceless }\OtherTok{\textless{}{-}} \FunctionTok{rnorm}\NormalTok{(}\AttributeTok{n =} \DecValTok{20000}\NormalTok{, }\AttributeTok{mean =} \DecValTok{170}\NormalTok{, }\AttributeTok{sd =} \DecValTok{44}\NormalTok{)}
\NormalTok{d.third.voiceless }\OtherTok{\textless{}{-}} \FunctionTok{rnorm}\NormalTok{(}\AttributeTok{n =} \DecValTok{20000}\NormalTok{, }\AttributeTok{mean =} \DecValTok{300}\NormalTok{, }\AttributeTok{sd =} \DecValTok{10}\NormalTok{)}

\NormalTok{d.voiced }\OtherTok{\textless{}{-}} \FunctionTok{data.frame}\NormalTok{(}\AttributeTok{category =} \StringTok{"voiced"}\NormalTok{,}\AttributeTok{VOT=}\NormalTok{ d.vot.voiced, }
                       \AttributeTok{VL =}\NormalTok{ d.vl.voiced, }\AttributeTok{Third  =}\NormalTok{ d.third.voiced)}

\NormalTok{d.voiceless }\OtherTok{\textless{}{-}} \FunctionTok{data.frame}\NormalTok{(}\AttributeTok{category =} \StringTok{"voiceless"}\NormalTok{, }\AttributeTok{VOT =}\NormalTok{ d.vot.voiceless, }
                          \AttributeTok{VL =}\NormalTok{ d.vl.voiceless, }\AttributeTok{Third =}\NormalTok{ d.third.voiceless)}

\NormalTok{d }\OtherTok{\textless{}{-}} \FunctionTok{rbind}\NormalTok{(d.voiced, d.voiceless)}

\CommentTok{\# sanity check}
\FunctionTok{ggplot}\NormalTok{(d, }\FunctionTok{aes}\NormalTok{(}\AttributeTok{x =}\NormalTok{ VOT)) }\SpecialCharTok{+}
  \FunctionTok{geom\_density}\NormalTok{()}
\end{Highlighting}
\end{Shaded}

\pandocbounded{\includegraphics[keepaspectratio]{toscano_sim2_files/toscano_sim2_files/figure-latex/unnamed-chunk-1-1.pdf}}

\section{2 build up the model architecture for a single cue
(e.g.~VOT)}\label{build-up-the-model-architecture-for-a-single-cue-e.g.-vot}

\begin{Shaded}
\begin{Highlighting}[]
\CommentTok{\# gap of understanding}
\NormalTok{mog.object }\OtherTok{\textless{}{-}} \ControlFlowTok{function}\NormalTok{(}\AttributeTok{K =} \DecValTok{20}\NormalTok{, }\AttributeTok{mu.m =} \DecValTok{25}\NormalTok{, }\AttributeTok{mu.sd =} \DecValTok{75}\NormalTok{, }\AttributeTok{sig =} \DecValTok{3}\NormalTok{, }\AttributeTok{seed =} \DecValTok{42}\NormalTok{)\{}
  \FunctionTok{list}\NormalTok{(}\AttributeTok{mu =} \FunctionTok{rnorm}\NormalTok{(K, mu.m, mu.sd),}
    \AttributeTok{sigma =} \FunctionTok{rep}\NormalTok{(sig, K),}
    \AttributeTok{phi =} \FunctionTok{rep}\NormalTok{(}\DecValTok{1} \SpecialCharTok{/}\NormalTok{ K, K))}
\NormalTok{\}}

\CommentTok{\# Get G(x) and M(x) according to equation 4 and 5}
\NormalTok{mog.G }\OtherTok{\textless{}{-}} \ControlFlowTok{function}\NormalTok{(mog.object, x) \{}
\NormalTok{  mog.object}\SpecialCharTok{$}\NormalTok{phi }\SpecialCharTok{*} \FunctionTok{dnorm}\NormalTok{(x, }\AttributeTok{mean =}\NormalTok{ mog.object}\SpecialCharTok{$}\NormalTok{mu, }\AttributeTok{sd =}\NormalTok{ mog.object}\SpecialCharTok{$}\NormalTok{sigma)}
\NormalTok{\}}
\NormalTok{mog.M }\OtherTok{\textless{}{-}} \ControlFlowTok{function}\NormalTok{(mog.object, x) }\FunctionTok{sum}\NormalTok{(}\FunctionTok{mog.G}\NormalTok{(mog.object, x))}

\NormalTok{mog.post }\OtherTok{\textless{}{-}} \ControlFlowTok{function}\NormalTok{(mog.object, x)\{}
\NormalTok{  G }\OtherTok{\textless{}{-}} \FunctionTok{mog.G}\NormalTok{(mog.object, x)}
\NormalTok{  M }\OtherTok{\textless{}{-}} \FunctionTok{mog.M}\NormalTok{(mog.object, x)}
\NormalTok{  G}\SpecialCharTok{/}\NormalTok{M}
\NormalTok{\}}

\CommentTok{\# learning rule according to A2{-}A5}
\CommentTok{\# eta is for learning rate}
\CommentTok{\# q.what is the use of sigma.min, sigma.max}
\NormalTok{mog.update }\OtherTok{\textless{}{-}} \ControlFlowTok{function}\NormalTok{(mog.object, x, }\AttributeTok{eta.mu =} \DecValTok{1}\NormalTok{, }\AttributeTok{eta.sigma =} \DecValTok{1}\NormalTok{, }\AttributeTok{eta.phi =} \FloatTok{0.001}\NormalTok{)\{}
\NormalTok{  r }\OtherTok{\textless{}{-}} \FunctionTok{mog.post}\NormalTok{(mog.object, x)}
\NormalTok{  mu.old }\OtherTok{\textless{}{-}}\NormalTok{ mog.object}\SpecialCharTok{$}\NormalTok{mu}
\NormalTok{  sg.old }\OtherTok{\textless{}{-}}\NormalTok{ mog.object}\SpecialCharTok{$}\NormalTok{sigma}

  \CommentTok{\# A3 }
\NormalTok{  mog.object}\SpecialCharTok{$}\NormalTok{mu }\OtherTok{\textless{}{-}}\NormalTok{ mu.old }\SpecialCharTok{+}\NormalTok{ eta.mu }\SpecialCharTok{*}\NormalTok{ r }\SpecialCharTok{*}\NormalTok{ (x }\SpecialCharTok{{-}}\NormalTok{ mu.old) }\SpecialCharTok{/}\NormalTok{ sg.old}\SpecialCharTok{\^{}}\DecValTok{2}
  \CommentTok{\# A4 }
\NormalTok{  sig.new }\OtherTok{\textless{}{-}}\NormalTok{ sg.old }\SpecialCharTok{+}\NormalTok{ eta.sigma }\SpecialCharTok{*}\NormalTok{ r }\SpecialCharTok{*}\NormalTok{ (sg.old}\SpecialCharTok{\^{}}\NormalTok{(}\SpecialCharTok{{-}}\DecValTok{3}\NormalTok{) }\SpecialCharTok{*}\NormalTok{ (x }\SpecialCharTok{{-}}\NormalTok{ mu.old)}\SpecialCharTok{\^{}}\DecValTok{2} \SpecialCharTok{{-}}\NormalTok{ sg.old}\SpecialCharTok{\^{}}\NormalTok{(}\SpecialCharTok{{-}}\DecValTok{1}\NormalTok{))}
\NormalTok{  mog.object}\SpecialCharTok{$}\NormalTok{sigma }\OtherTok{\textless{}{-}}\NormalTok{ sig.new}
  \CommentTok{\# ai suggests an extra sigma.max, sigma.min}
  \CommentTok{\#if (!is.na(sigma.max)) sg.new \textless{}{-} pmin(sg.new, sigma.max)}
  \CommentTok{\#mog.object$sigma \textless{}{-} pmax(sg.new, sigma.min)}

  \CommentTok{\# A2, winner{-}take{-}all, so that unnecessary categories don\textquotesingle{}t get updated}
\NormalTok{  winner }\OtherTok{\textless{}{-}} \FunctionTok{which.max}\NormalTok{(r)}
\NormalTok{  mog.object}\SpecialCharTok{$}\NormalTok{phi[winner] }\OtherTok{\textless{}{-}}\NormalTok{ mog.object}\SpecialCharTok{$}\NormalTok{phi[winner] }\SpecialCharTok{+}\NormalTok{ eta.phi }\SpecialCharTok{*}\NormalTok{ r[winner]}
\NormalTok{  mog.object}\SpecialCharTok{$}\NormalTok{phi }\OtherTok{\textless{}{-}} \FunctionTok{pmax}\NormalTok{(mog.object}\SpecialCharTok{$}\NormalTok{phi, }\DecValTok{0}\NormalTok{)}
\NormalTok{  mog.object}\SpecialCharTok{$}\NormalTok{phi }\OtherTok{\textless{}{-}}\NormalTok{ mog.object}\SpecialCharTok{$}\NormalTok{phi }\SpecialCharTok{/} \FunctionTok{sum}\NormalTok{(mog.object}\SpecialCharTok{$}\NormalTok{phi)}
\NormalTok{  mog.object}
\NormalTok{\}}

\CommentTok{\# calculate weight, equation 6, using outer product}
\NormalTok{mog.weight }\OtherTok{\textless{}{-}} \ControlFlowTok{function}\NormalTok{(mog.object)\{}
\NormalTok{  diffs }\OtherTok{\textless{}{-}} \FunctionTok{outer}\NormalTok{(mog.object}\SpecialCharTok{$}\NormalTok{mu, mog.object}\SpecialCharTok{$}\NormalTok{mu, }\ControlFlowTok{function}\NormalTok{(a, b) (a }\SpecialCharTok{{-}}\NormalTok{ b)}\SpecialCharTok{\^{}}\DecValTok{2}\NormalTok{)}
\NormalTok{  sigs  }\OtherTok{\textless{}{-}} \FunctionTok{outer}\NormalTok{(mog.object}\SpecialCharTok{$}\NormalTok{sigma, mog.object}\SpecialCharTok{$}\NormalTok{sigma, }\StringTok{"*"}\NormalTok{)}
\NormalTok{  phis  }\OtherTok{\textless{}{-}} \FunctionTok{outer}\NormalTok{(mog.object}\SpecialCharTok{$}\NormalTok{phi, mog.object}\SpecialCharTok{$}\NormalTok{phi, }\StringTok{"*"}\NormalTok{)}
  \FloatTok{0.5}\SpecialCharTok{*}\FunctionTok{sum}\NormalTok{(phis }\SpecialCharTok{*}\NormalTok{ diffs }\SpecialCharTok{/}\NormalTok{ sigs)}
\NormalTok{\}}

\CommentTok{\# sanity check on VOT only updating}
\NormalTok{mog }\OtherTok{\textless{}{-}} \FunctionTok{mog.object}\NormalTok{(}\AttributeTok{K =} \DecValTok{20}\NormalTok{, }\AttributeTok{mu.m =} \DecValTok{25}\NormalTok{, }\AttributeTok{mu.sd =} \DecValTok{75}\NormalTok{, }\AttributeTok{sig =} \DecValTok{3}\NormalTok{)}
\NormalTok{idx }\OtherTok{\textless{}{-}} \FunctionTok{sample}\NormalTok{(}\FunctionTok{nrow}\NormalTok{(d), }\DecValTok{90000}\NormalTok{, }\AttributeTok{replace =} \ConstantTok{TRUE}\NormalTok{) }\CommentTok{\# shuffle so voiced/voiceless aren\textquotesingle{}t blocked}
\ControlFlowTok{for}\NormalTok{ (i }\ControlFlowTok{in}\NormalTok{ idx) }
\NormalTok{  mog }\OtherTok{\textless{}{-}} \FunctionTok{mog.update}\NormalTok{(mog, d}\SpecialCharTok{$}\NormalTok{VOT[i])}

\FunctionTok{mog.weight}\NormalTok{(mog)}
\end{Highlighting}
\end{Shaded}

\begin{verbatim}
## [1] 10.44576
\end{verbatim}

\section{3 extend one cue to three cues (VOT, VL,
combined)}\label{extend-one-cue-to-three-cues-vot-vl-combined}

p446. Initial \(\mu\) values were randomly chosen from a distribution
with a mean of 25 and standard deviation of 75 for the VOT dimension and
a mean of 179 and standard deviation of 75 for the VL dimension. Initial
\(\sigma\) were set to 3 for the VOT dimiension and 10 for the VL
dimension. K was set to 20, and initial \(\phi\) were set to 1 ⁄ K.
Learning rates were set to 1 (\(\eta_\mu\)), 1 (\(\eta_\sigma\)), 0.001
(\(\eta_\phi\)), and 0.001 (\(\eta_p\)).2 The models were then tested on
a range of VOTs (0--40 ms in 5 ms steps) and two VLs (125 and 225 ms).

\begin{Shaded}
\begin{Highlighting}[]
\CommentTok{\# set initial parameters as somewhere between voiced and voiceless}
\NormalTok{vot.mog  }\OtherTok{\textless{}{-}} \FunctionTok{mog.object}\NormalTok{(}\AttributeTok{K =} \DecValTok{20}\NormalTok{, }\AttributeTok{mu.m =} \DecValTok{25}\NormalTok{,  }\AttributeTok{mu.sd =} \DecValTok{75}\NormalTok{, }\AttributeTok{sig =} \DecValTok{3}\NormalTok{)}
\NormalTok{vot.mog.raw  }\OtherTok{\textless{}{-}} \FunctionTok{mog.object}\NormalTok{(}\AttributeTok{K =} \DecValTok{20}\NormalTok{, }\AttributeTok{mu.m =} \DecValTok{25}\NormalTok{,  }\AttributeTok{mu.sd =} \DecValTok{75}\NormalTok{, }\AttributeTok{sig =} \DecValTok{3}\NormalTok{)}

\NormalTok{vl.mog }\OtherTok{\textless{}{-}} \FunctionTok{mog.object}\NormalTok{(}\AttributeTok{K =} \DecValTok{20}\NormalTok{, }\AttributeTok{mu.m =} \DecValTok{179}\NormalTok{, }\AttributeTok{mu.sd =} \DecValTok{75}\NormalTok{, }\AttributeTok{sig =} \DecValTok{10}\NormalTok{)}
\NormalTok{comb.mog }\OtherTok{\textless{}{-}} \FunctionTok{mog.object}\NormalTok{(}\AttributeTok{K =} \DecValTok{20}\NormalTok{, }\AttributeTok{mu.m =} \DecValTok{0}\NormalTok{,   }\AttributeTok{mu.sd =} \DecValTok{1}\NormalTok{,  }\AttributeTok{sig =} \FloatTok{0.2}\NormalTok{)}

\NormalTok{vot.gm }\OtherTok{\textless{}{-}} \FunctionTok{mean}\NormalTok{(d}\SpecialCharTok{$}\NormalTok{VOT); vot.gs }\OtherTok{\textless{}{-}} \FunctionTok{sd}\NormalTok{(d}\SpecialCharTok{$}\NormalTok{VOT)}
\NormalTok{vl.gm  }\OtherTok{\textless{}{-}} \FunctionTok{mean}\NormalTok{(d}\SpecialCharTok{$}\NormalTok{VL);  vl.gs  }\OtherTok{\textless{}{-}} \FunctionTok{sd}\NormalTok{(d}\SpecialCharTok{$}\NormalTok{VL)}
\NormalTok{w.vot }\OtherTok{\textless{}{-}} \FloatTok{0.5}
\NormalTok{w.vl  }\OtherTok{\textless{}{-}} \FloatTok{0.5}

\NormalTok{idx }\OtherTok{\textless{}{-}} \FunctionTok{sample}\NormalTok{(}\FunctionTok{nrow}\NormalTok{(d), }\DecValTok{5000}\NormalTok{, }\AttributeTok{replace =} \ConstantTok{TRUE}\NormalTok{)}
\ControlFlowTok{for}\NormalTok{ (i }\ControlFlowTok{in}\NormalTok{ idx) \{}
\NormalTok{  vot.val }\OtherTok{\textless{}{-}}\NormalTok{ d}\SpecialCharTok{$}\NormalTok{VOT[i]}
\NormalTok{  vl.val  }\OtherTok{\textless{}{-}}\NormalTok{ d}\SpecialCharTok{$}\NormalTok{VL[i]}
  \CommentTok{\# update the two cue MOGs from initial values}
\NormalTok{  vot.mog }\OtherTok{\textless{}{-}} \FunctionTok{mog.update}\NormalTok{(vot.mog, vot.val)}
\NormalTok{  vl.mog  }\OtherTok{\textless{}{-}} \FunctionTok{mog.update}\NormalTok{(vl.mog,  vl.val)}
  \CommentTok{\# normalize the weights so that they sum up to one}
\NormalTok{  w.vot.raw }\OtherTok{\textless{}{-}} \FunctionTok{mog.weight}\NormalTok{(vot.mog)}
\NormalTok{  w.vl.raw  }\OtherTok{\textless{}{-}} \FunctionTok{mog.weight}\NormalTok{(vl.mog)}
\NormalTok{  tot }\OtherTok{\textless{}{-}}\NormalTok{ w.vot.raw }\SpecialCharTok{+}\NormalTok{ w.vl.raw}
\NormalTok{  w.vot }\OtherTok{\textless{}{-}}\NormalTok{ w.vot.raw }\SpecialCharTok{/}\NormalTok{ tot;  w.vl }\OtherTok{\textless{}{-}}\NormalTok{ w.vl.raw }\SpecialCharTok{/}\NormalTok{ tot }
  \CommentTok{\# zscore the cue}
\NormalTok{  z.vot }\OtherTok{\textless{}{-}}\NormalTok{ (vot.val }\SpecialCharTok{{-}}\NormalTok{ vot.gm) }\SpecialCharTok{/}\NormalTok{ vot.gs}
\NormalTok{  z.vl  }\OtherTok{\textless{}{-}}\NormalTok{ (vl.val  }\SpecialCharTok{{-}}\NormalTok{ vl.gm ) }\SpecialCharTok{/}\NormalTok{ vl.gs}
\NormalTok{  x.comb }\OtherTok{\textless{}{-}}\NormalTok{ w.vot }\SpecialCharTok{*}\NormalTok{ z.vot }\SpecialCharTok{{-}}\NormalTok{ w.vl }\SpecialCharTok{*}\NormalTok{ z.vl}
  \CommentTok{\# update the combined ver as well}
\NormalTok{  comb.mog }\OtherTok{\textless{}{-}} \FunctionTok{mog.update}\NormalTok{(comb.mog, x.comb)}
\NormalTok{\}}
\end{Highlighting}
\end{Shaded}

\begin{verbatim}
## Warning in dnorm(x, mean = mog.object$mu, sd = mog.object$sigma): NaNs produced
## Warning in dnorm(x, mean = mog.object$mu, sd = mog.object$sigma): NaNs produced
\end{verbatim}

\begin{Shaded}
\begin{Highlighting}[]
\NormalTok{xs }\OtherTok{\textless{}{-}} \FunctionTok{seq}\NormalTok{(}\FunctionTok{min}\NormalTok{(d}\SpecialCharTok{$}\NormalTok{VOT) }\SpecialCharTok{{-}} \DecValTok{5}\NormalTok{, }\FunctionTok{max}\NormalTok{(d}\SpecialCharTok{$}\NormalTok{VOT) }\SpecialCharTok{+} \DecValTok{5}\NormalTok{, }\AttributeTok{length.out =} \DecValTok{400}\NormalTok{)}

\CommentTok{\# add a VOT axis}
\NormalTok{plot.df }\OtherTok{\textless{}{-}} \FunctionTok{rbind}\NormalTok{(}
  \FunctionTok{data.frame}\NormalTok{(}\AttributeTok{x =}\NormalTok{ xs, }\AttributeTok{version =} \StringTok{"raw"}\NormalTok{, }\AttributeTok{density =} \FunctionTok{sapply}\NormalTok{(xs, }\ControlFlowTok{function}\NormalTok{(x) }\FunctionTok{mog.M}\NormalTok{(vot.mog.raw, x))),}
  \FunctionTok{data.frame}\NormalTok{(}\AttributeTok{x =}\NormalTok{ xs, }\AttributeTok{version =} \StringTok{"updated"}\NormalTok{, }\AttributeTok{density =} \FunctionTok{sapply}\NormalTok{(xs, }\ControlFlowTok{function}\NormalTok{(x) }\FunctionTok{mog.M}\NormalTok{(vot.mog,x))))}

\FunctionTok{ggplot}\NormalTok{() }\SpecialCharTok{+}
  \FunctionTok{geom\_histogram}\NormalTok{(}\AttributeTok{data =} \FunctionTok{data.frame}\NormalTok{(}\AttributeTok{x =}\NormalTok{ d}\SpecialCharTok{$}\NormalTok{VOT), }\FunctionTok{aes}\NormalTok{(}\AttributeTok{x =}\NormalTok{ x, }\AttributeTok{y =} \FunctionTok{after\_stat}\NormalTok{(density)), }\AttributeTok{fill =} \StringTok{"grey85"}\NormalTok{,}\AttributeTok{bins =} \DecValTok{40}\NormalTok{,) }\SpecialCharTok{+}
  \FunctionTok{geom\_line}\NormalTok{(}\AttributeTok{data =}\NormalTok{ plot.df, }\FunctionTok{aes}\NormalTok{(}\AttributeTok{x =}\NormalTok{ x, }\AttributeTok{y =}\NormalTok{ density, }\AttributeTok{colour =}\NormalTok{ version)) }\SpecialCharTok{+}
  \FunctionTok{theme\_minimal}\NormalTok{()}
\end{Highlighting}
\end{Shaded}

\pandocbounded{\includegraphics[keepaspectratio]{toscano_sim2_files/toscano_sim2_files/figure-latex/unnamed-chunk-4-1.pdf}}

\section{The derivation of A2}\label{the-derivation-of-a2}

We want \(\frac{\partial \log M(x)}{\partial \mu_i}\).

Apply the log rule

\[
\frac{\partial \log M}{\partial \mu_i} = \frac{1}{M} \cdot \frac{\partial M}{\partial \mu_i}  = \frac{1}{M} \cdot \sum_j \frac{\partial G_j}{\partial \mu_i} 
\]

And as \(G_j\) only contains \(\mu_j\), not the other \(\mu\)'s. So
\(\frac{\partial G_j}{\partial \mu_i} = 0\) and only the \(j=i\) term
survives

\[
\frac{\partial G_j}{\partial \mu_i} = \frac{\partial G_i}{\partial \mu_i} = \phi_i \cdot \frac{1}{\sqrt{2\pi\sigma_i^2}} \cdot \frac{\partial}{\partial \mu_i}\left[ \exp!\left( -\frac{(x-\mu_i)^2}{2\sigma_i^2} \right) \right]
\]

Apply the exponential rule and the chain rule we get the following

\[
u(\mu_i) = -\frac{(x-\mu_i)^2}{2\sigma_i^2} \\[3pt]\frac{du}{d\mu_i} = -\frac{1}{2\sigma_i^2} \cdot \frac{d}{d\mu_i}\left[(x-\mu_i)^2\right]
\]

apply chain rule once again

\$\$ \frac{d}{d\mu_i}(x - \mu\_i)\^{}2 = 2(x - \mu\_i)
\cdot \frac{d}{d\mu_i}(x - \mu\_i) = 2(x - \mu\_i) \cdot (-1) = -2(x -
\mu\_i) \textbackslash{[}3pt{]} \frac{du}{d\mu_i} =
-\frac{1}{2\sigma_i^2} \cdot (-2)(x - \mu\_i) =
\frac{x - \mu_i}{\sigma_i^2}

\$\$

and plug it back

\$\$ \frac{\partial G_i}{\partial \mu_i} =\{\phi\_i
\cdot \frac{1}{\sqrt{2\pi\sigma_i^2}}
\cdot \exp!\left(-\frac{(x-\mu_i)^2}{2\sigma_i^2}\right)\}\cdot \frac{x - \mu_i}{\sigma_i^2}
\textbackslash{[}5pt{]}

\frac{\partial G_i}{\partial \mu_i} = G\_i(x)
\cdot \frac{x - \mu_i}{\sigma_i^2} \$\$

\end{document}
